\section{Beam-energy and rate hypotheses near 92 MeV}
\label{app:v505-energy}
Studies v5.5.0--v5.5.4 examine why the excesses extracted from different
campaigns need not agree under a common coupling. Two questions are distinct:
whether a feature moves in mass with beam energy, and whether its accepted
amplitude changes with energy after the usual radiative normalization.
A convenient mass hypothesis is
\begin{equation}
 m_y=m_*\left(\frac{E_y}{E_*}\right)^\alpha,
\end{equation}
with $\alpha=0$, $1/2$ and $1$ representing a fixed mass, square-root scaling
and linear scaling. These are comparisons among proposed shapes, not
established production laws. The later studies concentrate on a fixed
92 MeV region and relax the relative campaign rates.


With 92 MeV in 2016 as the reference, square-root mass scaling maps the
2015 and 2021 hypotheses to 62.338 and 117.317 MeV. Their signed fitted
roots in v5.5.0 are $+1.540$ and $-0.450$. Linear scaling instead maps them
to 42.240 and 149.600 MeV, with roots $-0.799$ and $+0.030$.
Neither prescription preserves a positive feature in all three campaigns.
These are fixed-hypothesis comparisons on different mass windows; the
study does not optimize a mass exponent or assign a calibrated preference
between those windows.

\widefig{1.0}{../figures/v550_beam_energy_scan_comparison.pdf}{The v5.5.0
scan comparison under fixed, square-root and linear beam-energy mass
scaling. The mapped responses test the proposed correspondence between
campaigns; their alignment is not evidence for an energy-dependent particle mass.}
{fig:v505-mass-scaling}

\subsection{A common experiment for the rate comparison}
For campaign $y$ and bin $j$, the model can be written
\begin{equation}
 \lambda_{yj}=b_{yj}+(L_y\theta_y)_j+
 a\,K_y(m_y,\sigma_y)\,R(E_y)\,t_{yj}(m_y,\sigma_y),
 \qquad a\geq0,
\end{equation}
where $t_{yj}$ is the bin-integrated Gaussian signal shape, $K_y$ is the
inherited yield conversion and each background nuisance vector has penalty
$\theta_y^T\theta_y/2$. The reference energy is $E_*=2.30$ GeV. The three
rate laws compared in v5.5.4 are
\begin{equation}
 R(E)=1,\qquad R(E)=(E/E_*)^\beta,\qquad
 R(E)=\exp[-k(E-E_*)].
 \label{eq:v505-rate-laws}
\end{equation}
The first is the common-coupling model. In the other two, $a$ is an
equivalent accepted-amplitude normalization at $E_*$; the fitted energy
factor has no demonstrated production or detector interpretation.
Only three beam energies constrain each two-parameter rate law.

The comparison holds the selected bins and GP constraints fixed across
models. The density conversion is evaluated at the 92 MeV MC reference;
shape variations change the bin-integrated template rather than silently
renormalizing the experiment. The power exponent is searched on $[-6,6]$.
Centroid studies subsequently allow a common mass over 90--94 MeV or separate
campaign centroids, and resolution studies compare the stated fixed and
constrained widths. Those changes introduce additional search freedom.

\begin{table}[htbp]\centering\small
\begin{tabular}{lrr}\toprule
Rate law & $\widehat a$ & Conditional 68\% profile slope interval\\\midrule
Common & $3.102\times10^{-6}$ & ---\\
Power & $8.516\times10^{-6}$ & $\beta=-2.987\;[-3.517,-2.470]$\\
Exponential & $9.817\times10^{-6}$ & $k=1.681\;[1.257,2.087]~\mathrm{GeV}^{-1}$\\
\bottomrule\end{tabular}
\caption{v5.5.4 fits at fixed 92 MeV and fully scaled widths. The common
normalization has a conditional 68\% log-Wald interval
$[2.066,4.657]\times10^{-6}$. These uncertainties condition on the
background and response model.}\label{tab:v505-rates}\end{table}

\widefig{1.0}{../figures/v554_rate_fits_with_uncertainties.pdf}{The v5.5.4
common, power-law and exponential rate fits. Shading gives approximate
68\% and 95\% pointwise log-Wald bands, including normalization--slope
covariance. It is not an uncertainty band on the validity of the energy law.}
{fig:v505-rates}

\subsection{Fit improvement and the cost of additional freedom}
At fixed 92 MeV and scaled widths, the common experiment gives raw likelihood
roots of 2.447 for common coupling, 4.180 for the power law and 4.146 for the
exponential. A root $\sqrt{Q_0}$ is not automatically a Gaussian significance
when a rate slope has been fitted. The v5.5.4 Gaussian reference includes
that slope search: the power-law reference evaluated at the Poisson
statistic gives $Z_G=3.809$. Allowing a common centroid to vary over
90--94 MeV gives a slightly larger root, 4.206, but the earlier matching
Gaussian reference gives $Z_G=3.528$. Better likelihood alone does not imply
stronger evidence once the additional hypotheses are counted.

The most flexible retained shape fit has raw root 4.598 and a broader
independent-amplitude reference-envelope value $Z_G=3.443$. That envelope
is neither the fitted model's own calibrated probability nor a bound over
all continuous shapes. These calculations use a predictive Gaussian
reference with $V_y=\operatorname{diag}(\widehat\lambda_{0,y})+L_yL_y^T$;
they do not calibrate direct Poisson tails, the original mass selection or
the full search. The note's resolution-based Sidak factor is not transferred
to these expanded model families.

\widefig{1.0}{../figures/v554_likelihood_and_local_reference_comparison.pdf}{Raw
likelihood roots and the corresponding conditional reference probabilities
for the retained rate and shape hypotheses. Triangles denoting broader
reference families must not be read as the tested energy model's own tail.}
{fig:v505-energy-reference}
\widefig{1.0}{../figures/v554_extraction_model_comparison.pdf}{The same observed
residuals under common coupling, an energy-dependent rate, and additional
shape freedom. The profiled background adjustment remains visible alongside
the signal contribution. These are alternative descriptions of the same
selected data.}{fig:v505-energy-extraction}

\subsection{Combining extracted amplitudes and count spectra}
The extracted-signal comparison checks whether a second representation of
the same data changes the combination. In the fixed Gaussian experiment,
the signed estimates and their covariance retain all amplitude information;
the extracted and joint amplitude log likelihoods agree to
$5\times10^{-14}$ in the saved numerical check. Independent campaigns
already enter through a product likelihood, so fitting their extracted
amplitudes supplies no additional independence factor.

With independent positive rates, the fixed Gaussian statistic is
\begin{equation}
 Q_G=\sum_{y=1}^3\max(z_y,0)^2,\qquad
 Q_G\sim\tfrac18\chi_0^2+\tfrac38\chi_1^2+
          \tfrac38\chi_2^2+\tfrac18\chi_3^2.
\end{equation}
Signed Stouffer and Fisher combinations test different alternatives. Choosing
among them after seeing their results would require another selection
correction. Without a common rate law, an upper bound on the total expected
signal rows, $T=\sum_yK_y\eps_y^2$, remains defined. The Gaussian-reference
construction used in v5.5.4 is
\begin{equation}
 U_{90}=\widehat T-s_T\Phi^{-1}
 \left[0.1\,\Phi(\widehat T/s_T)\right].
\end{equation}
At 92 MeV it gives 30,944 rows for common coupling and 36,523 for independent
rates. These are conditional known-covariance bounds, not calibrated
uncertainties on an absolute production rate.

The studies motivate checks of background modeling and campaign-dependent
response after unblinding. They do not establish a beam-energy law or
supply a calibrated acceptance correction. Version-by-version changes and
the standalone PDFs are indexed in the accompanying study archive.

\widefig{1.0}{../figures/v554_spectra_vs_extracted_significance.pdf}{Comparison
of simultaneous Poisson fits to count spectra with Gaussian fits to signed
extracted amplitudes. Their near agreement is a check of the retained
amplitude information. Alternative Stouffer and Fisher rules answer different
questions and do not add information to the joint fit.}{fig:v505-extracted}
